\subsection{Pure-candidate recovery does not identify a pure neural code}
\label{sec:composite-identification}

Recovering a generating member of a six-candidate library answers a closed-set
question: which candidate is favored if exactly one candidate generated the
signal? It does not establish that the underlying code is pure. In the frozen
$N=60$ stress analysis, the locked heuristic maximin allocation achieved a
minimum matched-pure recovery of $0.750$, yet an equal-coefficient mixture of
gain and update elicited a pure call in $0.800$ of replicates. The flexible
adequacy gate therefore did not provide composite-code protection. Separately,
the certified asymptotic design had worse minimum finite recovery than the
heuristic schedule ($0.665$ versus $0.750$). An optimization certificate for a
pairwise separation objective is not a finite-sample certificate of mechanism
identification.

The relevant geometry is the geometry of measured, nuisance-adjusted means.
For $y=HXa+Z\gamma+\varepsilon$, with sensor/sampling operator $H$, nuisance
matrix $Z$, and known Gaussian covariance $\Sigma$, let $W\Sigma W^{\mathsf T}=I$
and let $Q$ span the candidate columns after whitening and projection off $WZ$.
Then $z=Q^{\mathsf T}Wy$ has mean $Ba$, where $B=Q^{\mathsf T}WHX$ and
$r=\operatorname{rank}(B)$. Pure candidates are rays
$\mathcal P_j=\{a b_j:a\geq0\}$; a prespecified positive two-candidate class is
\begin{equation}
 \mathcal M_{jk}(\eta)=
 \{a[w b_j+(1-w)b_k]:a\geq0,\ \eta\leq w\leq1-\eta\}.
 \label{eq:bounded-composite-family}
\end{equation}
The coefficients refer to a fixed full-grid standardization, not fractions of
neurons or measured energy. These classes must be formed after the allowed
measurement and nuisance transformations; separation in event space alone
does not suffice.

Two limitations follow. First, a pure ray may intersect a composite cone even
when all pure rays differ, producing exactly identical observation laws.
Second, no informative finite-sample procedure can uniformly certify exact
purity against every arbitrarily small added component. For a pure mean $\mu$
and $\mu_\epsilon=(1-\epsilon)\mu+\epsilon\nu$, whitened Gaussian laws satisfy
\begin{equation}
 \operatorname{TV}(P_{\mu_\epsilon},P_\mu)
 \leq \tfrac12\epsilon\|\nu-\mu\|_2.
\end{equation}
If a rule certifies the pure code with probability $p$ under $\mu$, its
false-pure probability under $\mu_\epsilon$ is at least
$p-\epsilon\|\nu-\mu\|_2/2$. Uniform control at level $\alpha$ as
$\epsilon\downarrow0$ therefore forces $p\leq\alpha$. A scientifically
declared component floor, or an approximate-purity target, is necessary;
uniform power additionally requires a signal-strength floor.

\paragraph{A compatible-set alternative.}
With known covariance and a fixed design, the common confidence region
\begin{equation}
 \mathcal C_\alpha(z)=
 \{\mu:\|z-\mu\|_2^2\leq\chi^2_r(1-\alpha)\}
\end{equation}
can be intersected with every pure ray, bounded composite cone, and the null.
We retain every intersecting class. A bounded-family pure certificate is issued
only when exactly one pure ray remains, every composite class is excluded, and
the null is excluded. If the generating mean belongs to an allowed composite,
then coverage of that mean necessarily retains its composite class and
prevents certification. Hence false-pure probability is at most $\alpha$
over the entire declared continuous composite family, not merely a finite
grid of tested weights. This uses one simultaneous coverage event and does not
require knowing the generating pair. An empty set is reported as no
prespecified class compatible, not as support for the least-bad candidate.

\paragraph{The finite-information cost.}
A separately specified diagnostic reused the locked allocation and globally
standardized Table~\ref{tab:candidates} signals, with known unit Gaussian
noise, $\eta=0.25$, and $\alpha=0.05$. All fifteen pairs were evaluated at
weights $0.25$, $0.50$, and $0.75$, with $2{,}000$ replicates per generating
condition. A separate level-$0.05$ full-span residual check could remove
certificates but could not increase false-pure probability. The returned
within-span compatible set has at least $95\%$ true-class coverage; combining
its rejection with the separate residual gate provides only a conservative
$90\%$ joint non-rejection bound, not a $95\%$ overall adequacy guarantee.
At $N=60$ the
worst tested mixture false-pure rate was $0.0035$, but the minimum correct
pure-certificate rate was only $0.0060$. Repeating the same allocation with
independent idealized observations gave respective rates $0.0040$ and
$0.1375$ at $N=240$, and $0.0030$ and $0.7770$ at $N=960$. The analytic
guarantee, rather than these Monte Carlo rates, provides continuous-class error
control. These sample counts are not animal counts or executable protocol
recommendations.

The severe $N=60$ ambiguity has a geometric explanation: the nearest bounded
composite lies at squared distance $0.997$--$3.633$ from each of the six pure
means, below the common rank-six confidence radius squared $12.592$.
Thus even a confidence ball centered at the noise-free pure mean intersects
a composite class for every candidate. The result is a conditional cost of
honest mechanism exclusion, not a free improvement over the original recovery
rule. Its positive two-component, known-covariance, fixed-link assumptions
exclude arbitrary signed, higher-order, nonlinear and biological alternatives;
the link-function ambiguities identified above remain in force.
